\PassOptionsToPackage{dvipsnames, table}{xcolor}
\documentclass[dvipsnames, table, 11pt]{article}

\textwidth=7in
\textheight=9.5in
\hoffset=-1in
\voffset=-1in
\parskip=10pt
\parindent=0pt

\usepackage{style}

\begin{document}
     \title{Inferring critical transitions from timeseries}

     \author[*]{\textbf{PhD student:} Davide Papapicco}
     \author[*]{\authorcr \textbf{Supervision:} Graham Donovan}
     \author[*]{Lauren Smith}
     \author[$\dagger$]{Merryn Tawhai}

     \affil[*]{Department of Mathematics, Faculty of Science, University of Auckland, Auckland, New Zealand}
     \affil[$\dagger$]{Auckland Bioengineering Institute, University of Auckland, Auckland, New Zealand}
 
     \date{}

     \maketitle

     \begin{abstract}
                Critical transitions (or tipping points) are sudden, abrupt and often unforeseen departures of states of a dynamical system from a regime of stable equilibrium.
                In the past 20 years, tipping points have been an ever increasing studied subject in climate science and ecosystems given how disruptive these events can be.
                Mathematical characterisation of tipping points happened more recently with an tentative classification into three different mechanisms being proposed in 2012 by Ashwin et al.: bifurcation (B-tipping), noise (N-tipping) and rate (R-tipping) induced.
                Lately, more and more efforts from the community have shifted towards the robust and reliable prediction of these transitions in the form of early-warning signals (EWS).
                In this talk we will discuss one recent statistical method to derive an EWS from a large deviation principle using only timeseries data prior to a saddle-node type bifurcation.
     \end{abstract}
\end{document}
